\section*{\S~9a Ausbilder}
 
\label{sec:ausbilder}
\addcontentsline{toc}{section}{\S~9a Ausbilder}

\begin{enumerate}[label=(\arabic*)]
    \item Ausbilder für eine Kampfrolle werden durch die Leitung auf Empfehlung der bestehenden Ausbilder für diese Kampfrolle ernannt. Bestehen keine Ausbilder, so spricht der Stab eine Empfehlung aus. Dieser Empfehlung ist durch die Leitung immer Folge zu leisten. Eine Ablehnung bedarf einer Erklärung. Diese muss dem S1-Feldwebel, dem Personal-Feldwebel, dem UvD, dem GvD und den bestehenden Ausbildern vorgelegt werden, welche über die Gültigkeit dieser Erklärung abstimmen.
    \begin{enumerate}[label=(\Alph*)]
        \item Analog dazu kann die Leitung auch Personen vom Ausbilderposten entheben. Hierbei können auch der Stab, der S1-Feldwebel oder der Personal-Feldwebel eine solche Empfehlung aussprechen, selbst wenn bestehende Ausbilder vorhanden sind. Alle Empfehlungen, einen Ausbilder seines Postens zu entheben, sind durch die Leitung gründlich zu prüfen.
    \end{enumerate}
    \item Ausbilder kann jeder werden, der die Kampfrolle, die er ausbilden möchte, seit mindestens drei Monaten besitzt und sich mit dem Ausbildungsleitfaden/Handbuch vertraut gemacht hat. Pro Kampfrolle sollte es zwei, maximal vier Ausbilder geben. Ausbilder sollten aktive Mitglieder sein. Mindestens ein Ausbilder pro Kampfrolle muss ein aktives Mitglied sein, betrifft die Ausbilder der Sonder-Kampfrollen nicht.
    \begin{enumerate}[label=(\Alph*)]
        \item Eine Ausnahme hiervon ist möglich, wenn ein Mitglied eine Kampfrolle samt Ausbildungsleitfaden/Handbuch neu einführt oder den bestehenden Ausbildungsleitfaden/Handbuch in besonderem Maße erweitert oder überarbeitet. Eine solche Entscheidung obliegt dem Stab. Die bestehenden Ausbilder einer Kampfrolle sollten hierbei in beratender Funktion unterstützen.
        \item Die bestehenden Ausbilder können ein Mitglied, das Ausbilder werden möchte, auf die Eignung zur Ausbildung der Kampfrolle testen, bevor sie eine Empfehlung aussprechen. Wie sie dies tun, ist ihnen freigestellt.
        \item Für Sonder-Kampfrollen-Ausbilder kann die Leitung ggf. weitere Regelungen erlassen. Diese müssen klar und in geeigneter Weise kommuniziert werden.
    \end{enumerate}
    \item Ausbilder sind verpflichtet, ihren Ausbildungsleitfaden/Handbuch bei Bedarf anzupassen. Geschieht dies nicht, kann der betreffende Ausbilder seines Postens enthoben werden.
\end{enumerate}
